% ECONOMICS_PROBLEM: {"id": "I11-618", "title": "Economic allocation of mental-health care", "jel_code": "I11", "source_id": "OP-0058"}
% !TeX program = xelatex
\documentclass[11pt,a4paper]{article}
\usepackage[margin=23mm]{geometry}
\usepackage{fontspec}
\setmainfont{Latin Modern Roman}
\usepackage{amsmath,amssymb,mathtools}
\usepackage{xcolor,enumitem,booktabs,longtable,array,xurl,hyperref}
\definecolor{muted}{HTML}{53636C}
\hypersetup{colorlinks=true,urlcolor=blue,linkcolor=blue}
\setlength{\parindent}{0pt}
\setlength{\parskip}{5pt}
\newcommand{\R}{\mathbb R}
\newcommand{\E}{\mathbb E}
\newcommand{\N}{\mathbb N}
\newcommand{\ind}{\mathbf 1}
\newcommand{\Var}{\operatorname{Var}}
\newcommand{\Cov}{\operatorname{Cov}}
\newcommand{\supp}{\operatorname{supp}}
\DeclareMathOperator*{\argmin}{arg\,min}
\DeclareMathOperator*{\argmax}{arg\,max}
\newcommand{\entrymeta}[3]{{\sffamily\small\color{muted}\textbf{\texttt{#1}}\quad|\quad #2\par #3\par}}
\newcommand{\Problemref}[1]{\texttt{#1}}
\newcommand{\JELref}[2]{\texttt{#2} (JEL \texttt{#1})}
\begin{document}
\section*{Economic allocation of mental-health care}
\textbf{Problem ID:} \texttt{I11-618}\quad \textbf{JEL:} \texttt{I11}\par
% BEGIN ECONOMICS STATEMENT
\label{problem:OP-0058}
\label{atlas:prob:H8}
\noindent\textbf{Original atlas identifier:} H8; entry 58. \textbf{Field:} Health economics. \textbf{Original tractability label:} Easy.

\noindent\textbf{Classification:} Parameterized theoretical/empirical research agenda; not a certified unsolved theorem.

\subsection*{Definitions and assumptions}
Let patients $i=1,\ldots,n$ have pretreatment observable covariates $X_i$ and latent need $\theta_i$. Treatment $k\in\{0,1,\ldots,K\}$ includes no treatment ($k=0$), with a defined intensity, duration, and provider qualification. Potential outcomes $H_i(k)$, $Y_i(k)$, and $E_i(k)$ represent health, employment-related resources, and educational outcomes over a fixed horizon. Family/network spillovers require joint assignments $\mathbf k$ and replace these outcomes by $H_i(\mathbf k)$, etc. Specify a monetary welfare map $V_i$ rather than add health, earnings, and education in incompatible units. Costs $c_i(k)$ are real treatment and opportunity costs, and $b_k$ are provider-capacity limits. Fix welfare weights $w_i\geq0$ and budget $B\geq0$.

\subsection*{Formal research question}
\emph{Editorial formulation: the mathematical target below is not a theorem quoted from the references.}

For a measurable allocation rule $\pi$ based only on admissible pretreatment information, characterize
\[
 \sup_{\pi}\mathbb E\sum_iw_i\left[V_i(H_i(\pi),Y_i(\pi),E_i(\pi))-V_i(H_i(0),Y_i(0),E_i(0))\right]
\]
subject to $\mathbb E\sum_ic_i(\pi)\leq B$ and stated provider-capacity constraints. Define whether capacity applies in expectation or for every realized assignment. Identify heterogeneous intention-to-treat effects of randomized access, then distinguish them from treatment-received effects under explicit compliance assumptions. Incorporate treatment interactions and spillovers if relevant. Report an identified set of allocation values and $\varepsilon$-optimal rules, with sensitivity to latent severity, dropout, and the monetary valuation of health and education.

\subsection*{Interpretation and scope}
A ratio of average benefits to dollars need not maximize total welfare when treatments are indivisible, capacities bind, or beneficiaries overlap; the allocation problem is the rigorous target. Earnings can be a component of resource gain or a proxy for well-being, but should not be counted again through an education valuation already incorporating lifetime earnings. Diagnosis codes may respond to access, so stratification should use pretreatment information and an explicit measurement model. Waiting-list randomization identifies the specified offer effect, not every therapy's long-run causal return. Multiyear follow-up should include initial assignees and family outcomes, with missing-data assumptions stated. The original Easy classification describes the atlas author's judgment; identifying heterogeneous long-run welfare and spillovers is not thereby demonstrated easy.

\subsection*{Source audit and status}
The original sources are identifiable. [S1] and [S2] concern mental-health associations/outcomes, not randomized treatment returns; [S3] is a cost-benefit analysis. Publisher migration dates do not replace [S3]'s 2007 year. These works motivate the allocation question but do not certify a universally optimal or currently unresolved rule.

\subsection*{References}
\begin{enumerate}[label={[S\arabic*]},leftmargin=*,itemsep=0.6em]
\item Janet Currie and Mark Stabile (2006). Child Mental Health and Human Capital Accumulation: The Case of ADHD. Journal of Health Economics 25(6), 1094--1118. DOI: 10.1016/j.jhealeco.2006.03.001.
\url{https://www.sciencedirect.com/science/article/pii/S0167629606000282}
\newline\emph{Locator and support limits:} Abstract: ADHD symptoms and child outcomes; not randomized treatment returns.
\item Pinka Chatterji, Margarita Alegria, and David Takeuchi (2011). Psychiatric Disorders and Labor Market Outcomes: Evidence from the National Comorbidity Survey-Replication. Journal of Health Economics 30(5), 858--868. DOI: 10.1016/j.jhealeco.2011.06.006.
\url{https://www.sciencedirect.com/science/article/pii/S0167629611000774}
\newline\emph{Locator and support limits:} Abstract: psychiatric disorders and labor outcomes under selection assumptions; not direct randomized treatment effects.
\item Richard Layard, David Clark, Martin Knapp, and Guy Mayraz (2007). Cost-Benefit Analysis of Psychological Therapy. National Institute Economic Review 202(1), 90--98. DOI: 10.1177/0027950107086171.
\url{https://www.cambridge.org/core/journals/national-institute-economic-review/article/costbenefit-analysis-of-psychological-therapy/CBEC26AB2AB0ABC47D0C3F570E0FFA21}
\newline\emph{Locator and support limits:} Abstract and cost-benefit analysis: therapy expansion calculations. Later online migration dates do not change 2007 publication year.
\end{enumerate}

\subsection*{Common conventions: Source mathematical conventions}

Unless an entry states otherwise, agent and firm sets are finite. State and action spaces are measurable, strategies and outcome maps are measurable, probability laws are specified, and expectations exist for the targets under discussion. Infinite-horizon discount factors lie in $(0,1)$ and discounted objectives are finite. Norms, tolerances and complexity input sizes must be specified for computational comparisons. Constants or primitives that appear locally are inputs, not empirical facts asserted by this catalogue.

\textbf{Identification.} A statistical model $m\in\mathcal M$ implies an observable law $P_m$ and target $\theta(m)$. At law $P$, its identified set is
\[
\Theta(P;\mathcal M)=\{\theta(m):m\in\mathcal M,\ P_m=P\}.
\]
Point identification holds when this set is a singleton, or equivalently when $P_{m_1}=P_{m_2}$ implies $\theta(m_1)=\theta(m_2)$ for all compatible models. Sharp bounds mean bounds attained, or approached when the boundary is not attained, by compatible models. Writing this set is a definition, not a solution: the substantive task is to establish informative restrictions and derive or estimate the set. An unrestricted-confounding model may give only trivial bounds.

\textbf{Inference.} Identification, estimability and valid uncertainty quantification are separate questions. Fix $\alpha\in(0,1)$. A confidence procedure $C_n$ over a declared law class $\mathcal P$ has asymptotic uniform coverage at level $1-\alpha$ when
\[
\liminf_{n\to\infty}\inf_{P\in\mathcal P}\Pr_P\{\theta(P)\in C_n\}\geq1-\alpha.
\]
For set-identified targets the coverage event must be defined separately, for example coverage of the full identified set. If an entry uses identification-indexed models instead of uniquely indexed laws, the same coverage convention is applied over those models. Pointwise limits do not establish uniform validity, and asymptotic validity does not guarantee finite-sample precision. Sample sizes, dependence structures and weak-identification sequences are part of each application.

\textbf{Counterfactuals.} $Y_i(a)$ is a potential outcome under a specified intervention, including its timing, implementation and versions. It is not $Y_i$ conditional on voluntarily choosing $a$. A full intervention vector permits interference; shorthand scalar treatment notation does not silently impose no interference. Assignment, selection, attrition, anticipation and measurement must be specified. A claim about an instrument's local effect is not automatically a population policy effect.

\textbf{Economic equilibrium.} A counterfactual model includes best responses, constraints, feasibility, market clearing, state transitions and expectations consistent with the stated information. When several equilibria exist, either a declared selector is an input or the target is set-valued. Comparisons across policy regimes must use the stated selection convention. Reduced-form relationships are not presumed invariant after an equilibrium-changing intervention.

\textbf{Optimization and welfare.} Optimization is over the stated feasible set. If attainment is not established, $\sup$ or $\inf$ and $\varepsilon$-optimal choices replace an asserted maximizer or minimizer. For example, with value $V=\sup_{a\in\mathcal A}W(a)<\infty$, an $\varepsilon$-optimal set is $\{a\in\mathcal A:W(a)\geq V-\varepsilon\}$ for $\varepsilon>0$. Benefits, costs, risk and health or environmental outcomes can be aggregated only using declared units and conversion weights. Interpersonal utility weights, the treatment of fiscal transfers and foreign profits, discounting, and the behavioral welfare criterion are explicit inputs. A comparative-static derivative additionally requires existence and appropriate differentiability of the selected counterfactual.

\textbf{Sources.} Local labels [S1], [S2], and so forth restart in every question. Locators identify the relevant abstract, model, section or result at the level actually checked; a bibliographic or abstract-level verification is not represented as inspection of an entire proof. Working papers are distinguished from later publications and changing versions. Source summaries are paraphrased. All URL checks and editorial formulations refer to the audit date stated above.
% END ECONOMICS STATEMENT
\end{document}
